\subsection*{15. Make the Student's Thinking Visible}

A lecture is difficult to improve if the teacher cannot observe the minds supposedly being changed. Replace ``Does everyone understand?'' with a task that generates a trace: a 1-sentence prediction, a short computation, a diagram, a code commit, a ranked set of alternatives with reasoning, or a revision of a flawed argument.

A student may answer correctly by guessing, copying, or relying on a peer; a single trace is not conclusive. But carefully sampled individual work gives a teacher something to diagnose. In group settings, require a private 1st response before discussion; this prevents the confident participant from becoming the group's invisible answer engine. In technical lessons, preserve a small number of representative wrong answers with consent and appropriate anonymization. They are valuable input for the next instructional revision.

\subsection*{16. Active Learning Is About Cognition, Not Classroom Choreography}

``Active learning'' is sometimes reduced to breakout rooms, gamification, polls, and colorful worksheets. The meaningful part is the cognitive action: predicting, constructing, contrasting, arguing, testing, explaining, or revising. A 20-minute uninterrupted derivation can be active for a student who is following, checking, and predicting every step; a lively group discussion can be passive for the student who lets others think.

A major meta-analysis of 225 undergraduate STEM studies reported average advantages of active-learning approaches over traditional lecturing in examination performance and failure rates. That provides an important prior, not a universal instruction to ban exposition or expect the same effect in every course [S5]. An effective instructional design uses \textbf{direct explanation when it is efficient} and \textbf{meaningful participation when it improves learning}. The choice should be judged by the learner's later performance, not the teacher's ideology.

\subsection*{17. Questions Should Have Assigned Functions}

A diagnostic question tests prerequisites. A prediction question tests a mental model. A contrasting case clarifies category boundaries. An error critique develops detection. A decision question integrates trade-offs. A transfer question probes adaptability. A metacognitive question asks the learner how they know their answer should be trusted.

Bad questioning is a scatter of prompts inserted to keep a session lively. Good questioning is a deliberate sequence designed to reveal and alter a particular misconception. The instructor should write likely wrong answers before class and plan how to respond. ``What is the answer?'' is often less useful than ``What would change your answer?'' A facilitator should also know when silence is productive: individual thinking time can be more important than uninterrupted teacher speech.

\subsection*{18. Teach the Errors That Genuine Practice Produces}

A course that shows only correct worked examples teaches students a world in which important failures arrive already labeled. Reality rarely does. Engineers encounter unreliable logs, statisticians encounter bad sampling, consultants encounter conflicting stakeholder testimony, and researchers encounter findings that disappear when assumptions change.

For each important method, build an error library: a typical wrong approach, why it is attractive, the observation that disproves it, the correction, and the principle that travels to other cases. Ask learners to repair an intentionally faulty system. Later, ask them to invent a failure that would fool a superficial reviewer. This is especially valuable for technical professions in which competence consists partly of refusing an apparently successful but invalid result.

\subsection*{19. Productive Difficulty Must Be Accompanied by Repair}

Education often demands persistence. But difficulty without feedback can become a test of endurance rather than a mechanism of learning. If an assignment defeats half the cohort because package installation fails, the right response is infrastructure repair. If it defeats them because they cannot distinguish a causal effect from a correlation, the right response may be clearer contrasting cases and more practice. If they can identify the concept but cannot use it independently, the right response may be reduced scaffolding and an unseen task.

After a difficult activity, ask 3 questions: \textbf{What did the student try? What did the difficulty reveal? What new attempt will now be possible?} If the institution has no answer to the 3rd, it has designed a frustration event rather than an instructional loop.
